\chapter{A first design}
\label{ch:first}

This chapter takes the sample design of BAG from an empty folder to generated RTL. Then
it runs lint, simulation and formal on a smaller design. Run the commands in an empty
folder. The terms fabric, master, slave, bridge and clock-domain crossing (CDC) are
defined in the list in section~\ref{sec:wb-terms}.

\section{The sample workbook}

\cmd{bag template --sample} writes the sample workbook, \code{demo_ss}, and the file
header it names. \cmd{bag check} reads it and runs every check. \cmd{bag generate}
writes the design into the folder you give with \opt{-o}:

\lstinputlisting[style=console]{examples/console/first.txt}

The sample is a small subsystem with five fabrics:

\begin{itemize}
\item \code{main_xbar}, an AXI4 crossbar. Its masters are \code{cpu0}, \code{dma0} and
  \code{dbg0}, and the bridge \code{br_ahb2main} from the AHB-Lite bus. Its slaves are
  the memories \code{ddr0} and \code{sram0} and two bridges.
\item \code{cfg_bus}, an AXI4-Lite fabric behind the bridge \code{br_main2cfg}, with a
  second master \code{sec0}.
\item \code{periph_apb}, an APB4 fabric behind \code{br_main2per}.
\item \code{ahb_bus}, an AHB-Lite bus with the microcontroller \code{mcu0}, and
  \code{aon_apb}, an APB4 fabric behind it in its own clock domain.
\end{itemize}

Open \file{spec.xlsx} in Excel or LibreOffice to see how each sheet describes this.
Chapter~\ref{ch:workbook} explains every sheet.

\subsection{The two warnings}

The check gives no error and two warnings, so generation goes ahead. A
\term{warning} is a finding that does not stop BAG; an \term{error} does
(chapter~\ref{ch:checks}). Both warnings say the same thing for one master each:
\code{cpu0} has a 128-bit data path and \code{main_xbar} has 64 bits, so BAG puts a
width converter (\code{cpu0_dw_2_main_xbar}) between them. That converter cannot split a
WRAP burst, so it answers one with the error response SLVERR. The warning lists every
slave the master reaches through it. If your master sends WRAP bursts (a cache line
fill often does), give it the width of the fabric, or have it send INCR bursts.

\subsection{What was written}

\cmd{bag generate} wrote 136 files. The top level of \file{out} holds:

\begin{description}
\item[\file{rtl/}] the SystemVerilog: the wrapper \code{demo_ss_bus_top} and its
  sub-wrappers, the parameter package in \file{rtl/pkg/}, the BAG blocks in
  \file{rtl/inhouse/}, and the PULP blocks in \file{rtl/vendor/}.
\item[\file{demo_ss.f}] the file list, in compile order, for your own tools.
\item[\file{sw/}] the C header with the address map.
\item[\file{constraints/}] the clock and clock-domain-crossing constraints, in SDC and XDC.
\item[\file{doc/}] the memory map (Markdown and HTML), the port map, the tie-off report
  and the reset sequence.
\item[\file{tests/}] the generated cocotb tests and a copy of the IR they read.
\item[\file{formal/}] the SymbiYosys jobs and their property files.
\item[\file{bag_outputs.json}] the list of what was written, which \cmd{bag lint},
  \cmd{bag sim} and \cmd{bag formal} read.
\end{description}

Chapter~\ref{ch:outputs} describes every file. Run \cmd{bag generate} again and you get
the same bytes.

\section{A blank workbook}

Without \opt{--sample}, \cmd{bag template} writes a blank workbook. Checking it shows
what errors look like: each names the sheet, the row and the rule ID. With an error,
\cmd{bag check} exits with status 1, and \cmd{bag generate} writes nothing.

\lstinputlisting[style=console]{examples/console/first-blank.txt}

Fill in the sheets in the order of chapter~\ref{ch:workbook}, and run \cmd{bag check}
after each sheet.

\subsection{The least a design needs}

A design with one fabric needs these rows. Take $N$ masters and $M$ slaves, all on one
AXI4 crossbar in one clock domain:

\begin{itemize}
\item \sheet{Project}: the project name, the output directory and the five
  \code{Generate} settings, \code{Y} or \code{N}. These are the first errors above.
\item \sheet{Clocks} and \sheet{Resets}: one clock and one reset in one domain.
\item \sheet{Fabrics}: one row.
\item \sheet{Masters}: $N$ rows. \sheet{Slaves}: $M$ rows.
\item \sheet{Topology}: $N+M$ rows. Each master links to the fabric and the fabric
  links to each slave, with a \col{Port index} that counts from 0 on each side.
\item \sheet{Memory_Map}: at least one region for each slave, $M$ rows. A slave with
  no region is unreachable (\ruleid{MM-PATH-NO-REGION}, a warning).
\item \sheet{Address_Decode}: one row for each master, $N$ rows, with a miss response
  (\code{DECERR} on AXI4). A missing row is an error (\ruleid{DEC-MISSING-ROW}).
\item \sheet{Wrapper}: the top module name, the parameter package, and one row of
  \code{Exposed ports} for each master and slave, $N+M$ rows, all \code{Y}.
\end{itemize}

\sheet{Bridges}, \sheet{Access_Rules}, \sheet{Port_Options}, \sheet{User_Ports} and
\sheet{Hierarchy} may stay empty. \sheet{Hierarchy} keeps the one row for the top
wrapper that the template gives it. A workbook like this one, with two masters and two
slaves, checks with no finding and generates 60 files.

For the crossbar, \code{bag_axi_xbar} with \col{Latency mode} \code{NO_LATENCY} is a safe
default: it adds no spill register, and the example workbooks use it. Choose another mode
only to fix timing. Chapter~\ref{ch:workbook} explains every column.

The repository also holds three finished workbooks in its top-level \file{examples/}
directory (not the manual's \file{docs/manual/examples/}, and not part of a \code{pip install}).
Each checks with zero findings, and \cmd{bag flow} takes each from workbook to passing tests.
\file{apb_min.xlsx} is the smallest workbook that generates: one APB4 master and three slaves.
\file{mcu_ahb.xlsx} is a microcontroller bus: an AHB-Lite processor, an SRAM, and an APB bus of
three peripherals behind a \code{bag_ahb2apb} bridge. \file{soc_4x8.xlsx} is an AXI4 crossbar
with four masters and eight slaves, with an APB bus of six peripherals behind a
\code{bag_axi2apb} bridge; copy it as the start of a spec of your own.

\section{Lint, simulation and formal}

The rest of this chapter needs the EDA tools (section~\ref{sec:eda}). Proving the
sample design takes many hours of processor time, so the example uses a smaller design:
one APB4 fabric with one master and one slave. Its IR,
\file{apb4_min.ir.yaml}, is in \file{files/} of the examples, and chapter~\ref{ch:ir}
shows it. \cmd{bag generate} and \cmd{bag flow} accept an IR in place of a workbook.

\cmd{bag flow} runs every step in turn: check, generate, lint, compile, the generated
tests and, with \opt{--formal}, the formal jobs. It stops at the first step that fails.

\lstinputlisting[style=console]{examples/console/eda/first-flow.txt}

Each step prints a report and ends with a summary line. Lint must report no finding
that is not waived; the tests must all pass; every formal job must hold its assertions
and reach its covers. Chapter~\ref{ch:flows} explains each report.
